%% Chương 7: Kết luận
\clearpage
\phantomsection
\section*{CHƯƠNG 7. KẾT LUẬN VÀ HƯỚNG PHÁT TRIỂN}
\addcontentsline{toc}{section}{\numberline{}CHƯƠNG 7. KẾT LUẬN VÀ HƯỚNG PHÁT TRIỂN}
\setcounter{section}{7}
\setcounter{subsection}{0}
\setcounter{figure}{0}
\setcounter{table}{0}
\setcounter{equation}{0}

\subsection{Kết luận}

\subsubsection{Tóm tắt nghiên cứu}

Nghiên cứu này đề xuất, cài đặt, và đánh giá toàn diện một hệ thống mã hoá chọn lọc
DC H.264 lai ghép (hybrid), kết hợp ba tầng bảo mật: PLCM per-NALU keystream
generation, Arnold 2D Cat Map permutation, và cơ chế khuếch tán chuỗi XOR
(chained diffusion feedback). Thuật toán hybrid được so sánh với ba thuật toán
cơ sở (AES-128-CTR với fixed counter, DES-OFB với fixed IV, RC4 với cùng key
mỗi NALU) trên hai chiến lược (S1: toàn bộ VCL; S2: chỉ I-frame), sử dụng
10 chỉ số đánh giá từ 4 nhóm: bảo mật, chất lượng hình ảnh, tính đúng đắn, và hiệu năng.

\subsubsection{Kết quả chính}

Nghiên cứu thu được các kết quả chính sau:

\paragraph{Về tính đúng đắn:}
Tất cả 8 tổ hợp (4 thuật toán $\times$ 2 chiến lược) đạt BER = 0 và SHA-256 PASS,
xác nhận byte-exact decryptability trên toàn bộ file 219 MB. Đây là invariant cứng
của thiết kế XOR và được đảm bảo bởi tính xác định (determinism) của PLCM kết hợp
với tính khả nghịch (invertibility) của Arnold Cat Map.

\paragraph{Về bảo mật:}
Thuật toán hybrid nổi bật trên tất cả chỉ số bảo mật:
\begin{itemize}
  \item \textbf{NPCR S2 = 99.6903\% $> 99.6\%$:} Vượt ngưỡng Wu et al. \cite{wu_npcr},
        chứng minh hiệu quả của chained diffusion. Đây là thuật toán duy nhất
        có NPCR có ý nghĩa đánh giá (AES/DES/RC4: N/A do keystream cố định).
  \item \textbf{Entropy = 7.9996 $\approx$ 8.0:} Gần giá trị lý thuyết tuyệt đối,
        chứng minh keystream PLCM sinh ra bytes phân phối gần đều.
  \item \textbf{SSIM DC = 0.0031 $\approx$ 0:} Cấu trúc DC bị phá vỡ gần hoàn toàn.
  \item \textbf{TSSIM DC = 0.0071 $<$ baseline 0.0163:} Cấu trúc temporal bị phá vỡ.
        Các thuật toán baseline không đạt mục tiêu này.
\end{itemize}

\paragraph{Về hiệu năng:}
RC4 nhanh nhất (769 ns/NALU, S1), tiếp theo là AES (2.753 ns/NALU), DES (23.500 ns/NALU),
và Hybrid (29.245 ns/NALU). Hybrid chậm hơn RC4 $38\times$ nhưng DC encrypt time
tuyệt đối S1 = 556 ms --- chấp nhận được cho ứng dụng không real-time cứng.

Phát hiện pipeline convergence S2: khi chỉ 155 NALU được mã hoá, DC $< 1\%$ pipeline.
I/O đọc/ghi 219 MB chiếm $> 99\%$ tổng thời gian, làm tốc độ thuật toán không
còn ảnh hưởng đến thời gian tổng (Amdahl's Law).

\subsubsection{Đóng góp của nghiên cứu}

\begin{enumerate}
  \item \textbf{Thuật toán hybrid với per-NALU keystream:}
        Đề xuất kết hợp PLCM (keystream per-NALU via \texttt{nalu\_index}),
        Arnold 2D Cat Map (hoán vị), và chained diffusion feedback (khuếch tán)
        vào mã hoá DC H.264. Đây là sự kết hợp mới chưa được khảo sát trong
        các nghiên cứu mã hoá chọn lọc H.264 đã biết.

        Đặc biệt, việc dùng \texttt{nalu\_index} như ``tweak'' cho PLCM
        --- biến mỗi NALU có keystream độc lập --- giải quyết vấn đề keystream
        reuse căn bản của AES/DES/RC4 trong ngữ cảnh mã hoá chọn lọc.

  \item \textbf{Cơ chế chained diffusion:}
        Công thức $C[i] = (\text{permuted}[i] \oplus ks[i]) \oplus C[i-1]$
        tạo ra hiệu ứng avalanche trong 25 bytes DC: 1 bit thay đổi plaintext
        propagates đến tất cả bytes ciphertext phía sau. Kết quả NPCR S2 = $99.69\%$
        xác nhận tính hiệu quả thực nghiệm.

  \item \textbf{Tính byte-exact tuyệt đối (8/8 PASS):}
        Thiết kế đảm bảo giải mã hoàn toàn lossless cho mọi tổ hợp thuật toán
        và chiến lược. BER = 0 và SHA-256 PASS trên file 219 MB.

  \item \textbf{Kiến trúc IEncryption plugin-based (C++):}
        Interface chung cho phép thêm thuật toán mới với 0 thay đổi pipeline.
        Nền tảng cho triển khai phần cứng (FPGA) và nghiên cứu tiếp theo.

  \item \textbf{Phân tích và giải thích pipeline convergence S2:}
        Phát hiện và giải thích định lượng hiện tượng khi DC $< 1\%$ pipeline,
        tốc độ thuật toán không còn là bottleneck (Amdahl's Law trong ngữ cảnh
        mã hoá chọn lọc). Đây là nhận xét hữu ích cho thiết kế hệ thống thực tế.
\end{enumerate}

\subsection{Hạn chế của nghiên cứu}

\subsubsection{Hạn chế về keyspace}

\textbf{Vấn đề:} Tổng keyspace $2^{58}$ thấp hơn đáng kể so với ngưỡng $2^{128}$
khuyến nghị cho bảo mật hiện đại. Với GPU hiện đại có thể brute-force $10^{12}$
khóa/giây, $2^{58}$ bị crack trong vài ngày.

\textbf{Nguyên nhân:} PLCM sử dụng IEEE 754 double-precision (64-bit, 52-bit mantissa),
giới hạn keyspace hiệu dụng xuống $\approx 2^{53}$.
Arnold Cat Map đóng góp thêm $\approx 2^5$ nhưng không đủ bù.

\textbf{Đánh giá:} Với ngưỡng bảo mật hiện đại $2^{128}$, thuật toán hybrid
trong dạng hiện tại \textbf{không phù hợp} cho ứng dụng yêu cầu bảo mật cao.
Phù hợp hơn cho: (a) nghiên cứu thuật toán, (b) ứng dụng bảo vệ nội dung mức thấp
(perceptual security, không cần cryptographic security).

\subsubsection{UACI chưa đạt ngưỡng lý thuyết}

\textbf{Vấn đề:} UACI thực tế $\approx 32.5$--$32.7\%$, thấp hơn ngưỡng $33.46\%$
của Wu et al. Chênh lệch $\approx 0.7$--$1.0\%$.

\textbf{Nguyên nhân:} Fixed points của Arnold Cat Map $(P=3, Q=11, N=5)$:
byte 0 và byte 24 không bị hoán vị ($\pi(0) = 0$, $\pi(24) = 24$). Điều này
làm một số bit flip không được phân phối đều qua chained diffusion, kéo UACI xuống.

\textbf{Giải pháp đề xuất:} Thay Arnold fixed-parameter bằng key-dependent Arnold
--- tham số $(P, Q)$ được dẫn xuất từ khóa K theo từng NALU, loại bỏ fixed points
cố định. (Xem Hướng phát triển 7.3.2.)

\subsubsection{Hạn chế dataset}

\textbf{Vấn đề:} Toàn bộ thực nghiệm trên 1 file video duy nhất (\texttt{output.h264},
H.264 Baseline Profile). Không biết kết quả có tổng quát hoá không.

\textbf{Ảnh hưởng cụ thể:}
\begin{itemize}
  \item NPCR phụ thuộc phân phối DC bytes của video --- video khác có thể cho
        NPCR khác.
  \item SSIM/TSSIM phụ thuộc nội dung video (motion, texture complexity).
  \item Timing kết quả phụ thuộc hardware và OS scheduling.
\end{itemize}

\subsubsection{DC\_OFFSET hardcoded}

\textbf{Vấn đề:} DC offset = 19 bytes xác định bằng phân tích thực nghiệm
trên file output.h264 cụ thể (H.264 Baseline Profile). Với H.264 Main Profile
hoặc High Profile, cấu trúc slice header khác, DC bytes có thể ở offset khác.

\textbf{Rủi ro:} Nếu áp dụng pipeline này cho file H.264 Main/High Profile,
sẽ mã hoá sai vùng bytes, có thể gây crash decoder hoặc không đạt hiệu quả bảo mật.

\textbf{Giải pháp:} Parse đầy đủ H.264 bitstream (dùng h264bitstream) để tìm
offset DC đúng cho từng NALU theo profile --- thay vì hardcode.

\subsubsection{NPCR S1 chưa đạt ngưỡng}

NPCR S1 = 99.5260\% thấp hơn ngưỡng 99.6\% (0.074\% chênh). Điều này
do P/B-frames trong S1 có phân phối DC bytes phức tạp hơn I-frames,
ảnh hưởng đến hiệu ứng chained diffusion trung bình.
Cần nghiên cứu với nhiều video và key để xác định liệu đây là vấn đề hệ thống
hay chỉ xảy ra với video/key cụ thể này.

\subsection{Hướng phát triển}

\subsubsection{Tăng keyspace lên $\geq 2^{128}$}

\textbf{Đề xuất:} Thay PLCM double-precision ($2^{53}$) bằng PLCM multi-precision
128-bit fixed-point:

\begin{itemize}
  \item Biểu diễn $x_0$ và $p$ bằng số nguyên 128-bit (4 $\times$ uint32\_t hoặc
        2 $\times$ uint64\_t), thực hiện phép chia với precision đủ cao.
  \item Keyspace tăng lên $\approx 2^{127}$ (giảm 1 bit do ràng buộc $x_0 \in (0, 0.5)$).
  \item Chi phí tính toán: phép chia 128-bit chậm hơn double $\approx 3$--$5\times$,
        nhưng hybrid vẫn nhanh hơn nhiều so với full AES encryption.
\end{itemize}

Với giải pháp này, keyspace hybrid đạt $\approx 2^{127} + 2^5 \approx 2^{127}$,
vượt ngưỡng $2^{128}$ (thiếu 1 bit --- cần thêm 1 bit tham số nữa, ví dụ từ Arnold).

\subsubsection{Cải thiện UACI --- Key-dependent Arnold permutation}

\textbf{Đề xuất:} Thay Arnold fixed $(P=3, Q=11)$ bằng $(P_n, Q_n)$ được dẫn xuất
từ key K và nalu\_index $n$:

\begin{equation}
(P_n, Q_n) = \text{KeyDerive}(K, n, \{1,\ldots,N-1\}^2)
\end{equation}

Với $N = 5$: $P_n \in \{1, 2, 3, 4\}$, $Q_n \in \{1, 2, 3, 4\}$, 16 kết hợp.

\textbf{Lợi ích:}
\begin{itemize}
  \item Loại bỏ fixed points cố định (vị trí fixed point phụ thuộc vào $(P, Q)$,
        thay đổi mỗi NALU → không còn fixed point nhất quán).
  \item Tăng keyspace Arnold từ $2^5$ lên $2^5 \times N_{\text{nalu}} \approx 2^5 \times 2^{14} = 2^{19}$
        (với 19.037 NALU) --- tăng keyspace per-file.
  \item Dự báo UACI tăng gần $33.46\%$.
\end{itemize}

\subsubsection{Mở rộng dataset --- Nhiều video và profile}

\textbf{Đề xuất:} Kiểm tra trên ít nhất 10 video đa dạng:
\begin{itemize}
  \item \textbf{Resolution:} 480p (SD), 1080p (FHD), 4K (UHD).
  \item \textbf{Profile:} Baseline, Main, High (cần resolve vấn đề DC\_OFFSET).
  \item \textbf{Content type:} Animation (ít motion, DC đều), Sports (nhiều motion,
        DC phức tạp), Surveillance (ít thay đổi, DC ổn định), Live action.
  \item \textbf{Bitrate:} Low ($<$ 1 Mbps), Medium (5--10 Mbps), High ($>$ 20 Mbps).
\end{itemize}

Mục tiêu: xác lập khoảng NPCR = $[99.5\%, 99.7\%]$ cho hybrid trên nhiều video,
chứng minh tính tổng quát hoá của kết quả.

\subsubsection{Triển khai phần cứng FPGA}

\textbf{Động lực:} Tên đề tài nghiên cứu đề cập triển khai FPGA. Soft-core FPGA
(Xilinx Zynq hoặc Intel Cyclone) có thể đạt throughput cao hơn soft CPU nhờ:
\begin{itemize}
  \item Pipeline PLCM parallel: nhiều stage PLCM song song.
  \item Fixed-point arithmetic: 128-bit fixed-point nhanh hơn software.
  \item Streaming interface: nhận NALU bytes trực tiếp từ network, không qua file I/O.
\end{itemize}

\textbf{Mục tiêu hiệu năng:} Real-time encryption cho video 4K 30fps ($\approx$ 300 Mbps).
Ước tính: với 19.037 NALU/219 MB, tỷ lệ là $\approx$ 87 NALU/MB. Video 300 Mbps
cần xử lý $\approx$ 3.260 NALU/s. Với FPGA 100 MHz pipeline, cần $< 30.000$
clock cycles/NALU cho hybrid.

\subsubsection{Mở rộng sang H.265/HEVC và AV1}

\textbf{H.265/HEVC:} Thay macroblock 16$\times$16 bằng CTU (Coding Tree Unit)
64$\times$64, làm thay đổi cách DC coefficients được tổ chức trong bitstream.
Cần phân tích lại để tìm DC\_OFFSET tương đương trong H.265.

\textbf{AV1:} Open-source codec do Alliance for Open Media phát triển,
đang ngày càng phổ biến (YouTube, Netflix). Sử dụng superblock 128$\times$128
và entropy coding khác hoàn toàn. Nghiên cứu DC selective encryption cho AV1
là hướng nghiên cứu mới.

\subsubsection{Phân tích bảo mật hình thức}

\textbf{Đề xuất:} Phân tích toán học chặt chẽ các tính chất bảo mật của hybrid:
\begin{itemize}
  \item \textbf{IND-CPA:} Chứng minh (hoặc bác bỏ) tính IND-CPA của hybrid
        với keyspace $2^{128}$ (sau khi nâng cấp PLCM).
  \item \textbf{IND-CCA:} Phân tích khả năng tấn công chosen-ciphertext.
  \item \textbf{Chosen-IV attack:} Phân tích nếu kẻ tấn công kiểm soát được
        \texttt{nalu\_index} (ví dụ trong video streaming theo yêu cầu).
  \item \textbf{Differential cryptanalysis:} Phân tích mức độ an toàn của
        chained diffusion trước tấn công vi sai (differential attack).
\end{itemize}

\subsubsection{Kết hợp với kỹ thuật bảo mật luồng (Stream Security)}

\textbf{Đề xuất:} Kết hợp DC selective encryption với HMAC (Hash-based Message
Authentication Code) để đảm bảo tính toàn vẹn (integrity) của file mã hoá.
Hiện tại, kẻ tấn công có thể sửa DC bytes encrypted mà không bị phát hiện ---
hệ thống chỉ có confidentiality, không có integrity.

Giải pháp: thêm SHA-256 HMAC trên toàn bộ file encrypted như metadata ngoài băng tần
(out-of-band) để phát hiện tampering.

\subsection{Tổng kết}

Nghiên cứu này chứng minh rằng mã hoá chọn lọc DC H.264 với thuật toán hybrid
(PLCM + Arnold + Chained Diffusion) là khả thi và hiệu quả cho perceptual security.
Thuật toán đạt NPCR $> 99.6\%$ (S2), Entropy $\approx 8.0$, và byte-exact decryptability
trên file H.264 thực tế 219 MB, 19.348 NALU.

Hạn chế chính là keyspace $2^{58}$ chưa đủ cho cryptographic security, được thừa nhận
như hướng nghiên cứu tương lai. Đối với các ứng dụng bảo vệ nội dung mức perceptual
(giám sát an ninh nhẹ, watermarking, content protection không yêu cầu FIPS compliance),
thuật toán hybrid trong dạng hiện tại đạt hiệu quả tốt với overhead chấp nhận được.

Kiến trúc IEncryption plugin-based và pipeline C++ đầy đủ tạo nền tảng vững chắc
cho triển khai thực tế và nghiên cứu mở rộng trong tương lai, đặc biệt là hướng
FPGA hardware implementation cho video 4K real-time.
